En la tabla comparamos los dos estilos. SOAP exige un WSDL \cite{w3c_wsdl11}, usa XML y avisa los errores con un Fault; REST puede describirse con OpenAPI \cite{openapi31}, usa casi siempre JSON y responde con códigos HTTP \cite{rfc9110}. Con REST se puede usar la caché de HTTP \cite{rfc9111}. En nuestro PFC la misma consulta pesa 252 bytes por REST y 703 por SOAP.

Nuestra postura: SOAP conviene cuando muchas organizaciones deben cumplir un contrato fijo, como los emisores de facturas ante el SRI, y REST conviene para la web, las apps y cuando importa el tamaño. Pero REST mal hecho también cuesta: solo el 0,8\,\% de 500 API públicas cumple todos sus principios \cite{neumann2021}, romper las reglas de diseño la vuelve más difícil de entender \cite{bogner2023}, y cambiar una API sin romper a sus clientes exige versiones \cite{lercher2024}.
